\documentclass[12pt]{article}
\usepackage[utf8]{inputenc}
\usepackage[T1]{fontenc}
\usepackage[a4paper,margin=3cm]{geometry}
\usepackage{minted}
\usepackage{xcolor}
\usepackage{csquotes}

\setlength{\parindent}{0em}
\setlength{\parskip}{\medskipamount}

\title{Einführung in Manim}
\author{}
\date{}

\begin{document}
	
	\maketitle
	
	\section*{Einleitung}
	
	Manim bietet die Möglichkeit, mathematische Zusammenhänge zu animieren. Das wollen wir uns im nächsten Kapitel anschauen.
	
	\section{Wie sieht ein Manim-Programm aus}
	
	Beispiel:
	
	\begin{minted}[bgcolor=gray!10]{python}
from manim import * # Import von Manim

class Example(Scene):
  def construct(self):
    dot = Dot(color=BLUE)
    square = Square(color=GREEN, fill_opacity=0.5)
    
    dot.next_to(square, DOWN)
    self.add(square, dot)
	\end{minted}
	
	Wenn wir Manim-Programme mit Jupyter oder Colab rendern wollen, müssen wir in die erste Zeile der Zelle folgendes schreiben:
	
\begin{minted}[bgcolor=gray!10]{python}
%%manim -qm <Classname>
	\end{minted}
	
	wobei \verb|<Classname>| den Namen der jeweiligen Klasse beschreibt.
	
	\section{Mobjects}
	
	\begin{minted}[bgcolor=gray!10]{python}
from manim import *

class Shapes(Scene):
  def construct(self):
    circle = Circle()
    square = Square()
    triangle = Triangle()

    circle.shift(LEFT)
    square.shift(UP)
    triangle.shift(RIGHT)

    self.add(circle, square, triangle)
	\end{minted}
	
	Mögliche Mobjects sind u.A.
	\begin{itemize}
		\item \verb|Polygon|,
		\item \verb|Rectangle|,
		\item \verb*|RegularPolygon|,
		\item \verb*|Square|
		\item \verb*|Star|,
		\item \verb*|Triangle|,
		\item \verb*|Line|,
		\item \verb*|Angle| bzw. \verb*|RightAngle|,
		\item \verb*|Dot|,
		\item \verb*|Tex| bzw. \verb*|MathTex|
	\end{itemize}
	
	Der Koordinatenursprung wird mit \verb*|ORIGIN| bezeichnet. Die Richtungen sind \verb*|UP|, \verb*|DOWN|, \verb*|LEFT| und \verb*|RIGHT|.
	
	Methoden wie \verb|.shift()| können auch zusammengefügt werden. So haben folgende beiden Programmausschnitte die gleiche Funktion:
	
	\begin{minted}[bgcolor=gray!10]{python}
square = Square()
square.shift(DOWN)

square = Square().shift(DOWN)
	\end{minted}
	
	Andere Möglichkeiten zur Positionierung sind folgende:
	
	\begin{minted}[bgcolor=gray!10]{python}
from manim import *
	
class Shapes(Scene):
  def construct(self):
    circle = Circle()
	square = Square()
	triangle = Triangle()
	
	circle.move_to(2 * LEFT)
	square.next_to(circle, LEFT)
	triangle.align_to(circle, RIGHT)
	
	self.add(circle, square, triangle)
\end{minted}

	\section{Styling}
	
\begin{minted}[bgcolor=gray!10]{python}
from manim import *

class MobjectStyling(Scene):
  def construct(self):
    circle = Circle().shift(LEFT)
    square = Square().shift(UP)
    triangle = Triangle().shift(RIGHT)

    circle.set_stroke(color=GREEN, width=20)
    square.set_fill(YELLOW, opacity=1.0)
    triangle.set_fill(PINK, opacity=0.5)

    self.add(circle, square, triangle)
\end{minted}

Die Reihenfolge der Mobjects kann man durch die Reihenfolge in \verb*|.add| ändern: \mintinline{python}|self.add(circle, square, triangle)| vs. \mintinline{python}|self.add(triangle, circle, square)|

	\section{Animationen}
	
\begin{minted}[bgcolor=gray!10]{python}
from manim import *

class SomeAnimations(Scene):
  def construct(self):
    square = Square()

    self.play(FadeIn(square))

    self.play(Rotate(square, PI/4))

    self.play(FadeOut(square))

    self.wait(1)
\end{minted}

Wir können aber auch über die Methode \verb*|.animate| auf das Mobject zugreifen:

\begin{minted}[bgcolor=gray!10]{python}
from manim import *
	
class SomeAnimations(Scene):
  def construct(self):
	square = Square()
	
	self.add(square)
	
	self.play(square.animate.set_fill(WHITE))
	
	self.play(square.animate.shift(UP).rotate(PI/4), run_time = 3)
	
	self.play(FadeOut(square))
	
	self.wait(1)
\end{minted}

\begin{minted}[bgcolor=gray!10]{python}
from manim import *
	
class SomeAnimations(Scene):
  def construct(self):
	tex = MathTex(r'\frac{10}{x}')
	
	self.add(tex)
	
	self.play(tex.animate.shift(UP).rotate(PI/4), run_time = 3)
	
	self.play(FadeOut(tex))
	
	self.wait(1)
\end{minted}
	
\end{document}
